%% parts/p07/calculi-and-derivability.tex --- Phase-4 P7.1 migration
%% Migrated source: v1.1 Part VIII, Classical and Intuitionistic Sequent
%% Calculi, lines 11357--11710. Depth: D-P8-001/002; apparatus A-070/071.
%% Rules live in registry/calculi; proof material lives in annex/p07/.

\chapter{Calculi and Derivability}
\label{sec:p07-calculi-and-derivability}
\chaptermark{Calculi and Derivability}

Structural proof theory studies derivations as formal objects.  A statement
about admissibility, normalization, or proof translation is meaningful only
after the language, sequent shape, context discipline, and primitive rules
have been fixed.  The exact rule surfaces used here are
\hyperref[sheet:lk]{\texttt{sheet:lk}},
\hyperref[sheet:lk-defneg]{\texttt{sheet:lk-defneg}},
\hyperref[sheet:lk-eq]{\texttt{sheet:lk-eq}}, and
\hyperref[sheet:lj]{\texttt{sheet:lj}}.  The defined-negation presentation
is the one inherited from v1.1 Part~VIII; the canonical and equality
presentations remain distinct.  See
\cite{Gentzen1934a,Gentzen1935b,TroelstraSchwichtenberg2000,
NegriVonPlato2001}.

\section{Sequents, contexts, and calculus identities}

\begin{hbconvention}{def:p07-contexts}{Contexts and formula metavariables}
In Chapters~\ref{sec:p07-calculi-and-derivability} and
\ref{sec:p07-cut-elimination-and-interpolation},
\(\varphi,\psi,\chi\) range over
formulas and \(\Gamma,\Delta,\Sigma,\Pi\) over finite multisets of formulas.
Commas denote multiset union.  Exchange is built into this representation;
weakening and contraction are not.  Thus \(\Gamma,\varphi\) adds one
occurrence of \(\varphi\), and \(\operatorname{supp}(\Gamma)\) forgets
multiplicity.
\end{hbconvention}

\begin{hbdefinition}{def:p07-classical-intuitionistic-sequents}{Classical and intuitionistic sequents}
A \emph{classical sequent} is a syntactic expression
\(\Gamma\Rightarrow\Delta\), read semantically as
\(\bigwedge\Gamma\to\bigvee\Delta\), with
\(\bigwedge\varnothing=\top\) and \(\bigvee\varnothing=\bot\).  An
\emph{intuitionistic sequent} is a syntactic expression
\(\Gamma\Rightarrow\varphi\) having exactly one succedent formula.  The latter
restriction is part of the calculus identity, not an abbreviation for a
classical sequent.  Neither kind of sequent is itself an object-language
conditional or a semantic consequence assertion.
\index{sequent!classical}
\index{sequent!intuitionistic}
\end{hbdefinition}

\begin{hbdefinition}{def:p07-lkd-derivation}{LKd and LJ derivations}
Write
\[
 \Gamma\derives_{\mathsf{LK}_{\mathrm d}}\Delta
\]
when a finite rooted tree with endsequent
\(\Gamma\Rightarrow\Delta\) is generated by the registered rules of
\(\mathsf{LK}_{\mathrm d}\).  Write
\[
 \Gamma\derives_{\mathsf{LJ}}\varphi
\]
analogously for the endsequent
\(\Gamma\Rightarrow\varphi\) on
\hyperref[sheet:lj]{\texttt{sheet:lj}}.  A leaf is an instance of a
zero-premise rule, and every other node is the conclusion of exactly one
displayed rule instance whose premises are its immediate predecessors.  A
derivation is \emph{cut-free} iff it contains no Cut instance.

The subscript \(\mathrm d\) is necessary because
\(\mathsf{LK}_{\mathrm d}\) has defined negation, uncontexted identity, and a
split-context implication-left rule, whereas canonical \(\mathsf{LK}\) has
primitive negation and contexted identity.  It names a calculus, not a new
logic.  The translation below is the later use of this distinction; no
equality rule belongs to either \(\mathsf{LK}_{\mathrm d}\) or
\(\mathsf{LJ}\).
\index{LKd@\(\mathsf{LK}_{\mathrm d}\)}
\index{derivation!sequent calculus}
\end{hbdefinition}

\begin{hbtheorem}{thm:p07-lkd-soundness}{sketched=annex:p7.calculi.1}{Soundness of the defined-negation sequent calculus}
Let \(\Gamma,\Delta\) be finite contexts.  If
\(\Gamma\derives_{\mathsf{LK}_{\mathrm d}}\Delta\), then every
classical structure and assignment satisfying every formula occurrence in
\(\Gamma\) satisfies at least one formula occurrence in \(\Delta\).  For an
empty succedent, the conclusion says that no structure and assignment
satisfies every antecedent formula.
\end{hbtheorem}

Semantic soundness of \(\mathsf{LJ}\) is not inferred from this classical
result.  Its Kripke semantics and adequacy belong to the Part-8
intuitionistic-logic entry.

\begin{hbwarning}{def:p07-sequent-levels}{Three levels that typography does not identify}
The sequent \(\Gamma\Rightarrow\Delta\), the derivability assertion
\(\Gamma\derives_{\mathsf{LK}}\Delta\), and the assertion that every
structure and assignment satisfies the sequent are respectively a syntactic
object, a metalinguistic claim about a finite tree, and a semantic claim.
Soundness and completeness connect these levels; the arrow and turnstile do
not identify them.
\end{hbwarning}

\section{Primitive, derived, and admissible rules}

\begin{hbdefinition}{def:p07-rule-standing}{Primitive, derived, and admissible rules}
Relative to one fixed calculus, a \emph{primitive rule} occurs in its
definition.  A rule is \emph{derived} when one uniform derivation takes
arbitrary instances of its premises as open leaves to the corresponding
conclusion.  A rule is \emph{admissible} when derivability of all premises of
each substitution instance implies derivability of its conclusion.  Every
derived rule is admissible, but admissibility alone supplies no uniform
derivation with the premises as open leaves.  These standings are
nonredundant: Cut is primitive on the active sheets, while
Chapter~\ref{sec:p07-cut-elimination-and-interpolation} only sketches its
admissibility in their cut-free fragments.
\index{rule!primitive}
\index{rule!derived}
\index{rule!admissible}
\end{hbdefinition}

\begin{hbdefinition}{def:p07-principal-formula}{Principal and active formula occurrences}
In a logical-rule instance, the \emph{principal formula} is the compound
occurrence displayed in the conclusion and introduced by that instance; the
displayed premise occurrences from which it is formed are its
\emph{active formulas}; every other occurrence is context.  In weakening or
contraction, the displayed weakened or contracted occurrence is principal.
A zero-premise rule has no principal formula.  The distinction is needed by
the cut-reduction cases, which branch on whether the mix formula is principal
in either immediate subderivation.
\index{principal formula}
\index{active formula}
\end{hbdefinition}

The citation-only comparison rows below use Girard's IMLL presentation
\cite{Girard1987} and the basic ordered Lambek presentation surveyed in
\cite{GalatosJipsenKowalskiOno2007}; neither citation supplies a P7 registry
sheet.

\begin{table}[!htbp]
\centering
\begingroup\scriptsize
\setlength{\tabcolsep}{3pt}
\begin{tabular}{@{}p{.18\textwidth}p{.12\textwidth}p{.12\textwidth}
  p{.12\textwidth}p{.12\textwidth}p{.25\textwidth}@{}}
\toprule
Exact calculus & Exchange & Weakening & Contraction & Cut displayed? &
Admissibility evidence / scope \\
\midrule
Canonical \(\mathsf{LK}\), \texttt{sheet:lk} & Built into both multiset
contexts & Primitive left and right & Primitive left and right & Primitive &
Not established by the Part-VIII reduction, whose logical and identity rules
are those of \(\mathsf{LK}_{\mathrm d}\). \\
\(\mathsf{LK}_{\mathrm d}\), \texttt{sheet:lk-defneg} & Built into both
multiset contexts & Primitive left and right & Primitive left and right &
Primitive & Sketched in Theorem~\ref{thm:p07-cut-elimination}; equality and
primitive negation are outside its scope. \\
\(\mathsf{LJ}\), \texttt{sheet:lj} & Built into the antecedent multiset &
Primitive on the left; no right weakening in the exact-single-succedent
syntax & Primitive on the left; no right contraction & Primitive & Sketched
in Theorem~\ref{thm:p07-cut-elimination}. \\
Girard's commutative IMLL sequent presentation & Built into multiset context &
Absent & Absent & Displayed in the cited presentation & Cut elimination is
citation-only here; this row neither displays its rules nor implies a P7
registry sheet. \\
Associative Lambek calculus \(\mathsf L\) with ordered antecedents & Absent &
Absent & Absent & Displayed in the selected basic sequent presentation &
Cut admissibility is citation-only; the ordered context, rather than a
multiset, blocks exchange. \\
\bottomrule
\end{tabular}
\endgroup
\caption{Structural rules by exact presentation (apparatus A-070).  Primitive
Cut and admissibility are separate columns.  Relevance logics have no
family-wide structural profile: contraction and mingle vary by system and
presentation.  Their concrete Routley--Meyer/system-specific treatment is
Part-8-owned and is omitted rather than genericized.  The ledger fixes no
source-exact affine sheet, so no affine row is fabricated.}
\label{tab:p07-structural-rules}
\end{table}

\hbstub{A-070: the relevance and affine comparison rows remain open because
no concrete source-exact sequent presentation has been selected.  Part~8 owns
the relevance-system selection; a named commutative affine calculus and its
exact Cut evidence must be fixed before its structural profile can be entered.}

\section{Completeness and presentation translations}

\begin{hbmethod}{def:p07-lk-completeness-lookup}{Finite-sequent completeness lookup}
For the propositional restriction of canonical
\hyperref[sheet:lk]{\texttt{sheet:lk}}, finite-multiset
\(\Gamma\Rightarrow\Delta\) is derivable iff every valuation satisfying all
members of \(\operatorname{supp}(\Gamma)\) satisfies some member of
\(\operatorname{supp}(\Delta)\); empty contexts are included.  This is the
sketched Theorem~\ref{thm:prop-lk-completeness}.  For equality-free
first-order canonical \(\mathsf{LK}\), the same finite-sequent equivalence
over all structures and assignments holds assuming the Axiom of Choice; this
is the sketched Theorem~\ref{thm:fol-lk-completeness}.  Both results use
the finite-disjunction reduction to single-conclusion consequence and
presently use Cut.  The cut-elimination sketch below is about
\(\mathsf{LK}_{\mathrm d}\), so it does not establish cut-free completeness
for canonical \(\mathsf{LK}\).
\end{hbmethod}

This lookup creates no duplicate theorem or annex item at a different
evidence level.  Propositional and first-order soundness remain the proved
Theorems~\ref{thm:prop-lk-soundness} and
\ref{thm:fol-lk-soundness}.

\begin{hbdefinition}{def:p07-defined-negation-translation}{Translation from canonical LK to the defined-negation language}
For equality-free formulas, define \(d\) recursively by
\[
 d(\bot)=\bot,\qquad d(\neg\varphi)=d(\varphi)\to\bot,
\]
and let \(d\) commute with every other propositional connective and with
\(\forall,\exists\), preserving variables and terms.  Apply \(d\)
occurrencewise to contexts.  This is a metalinguistic formula translation,
not a connective or an assertion of formula identity.  It is needed because
canonical \(\mathsf{LK}\) has primitive negation while
\(\mathsf{LK}_{\mathrm d}\) does not; the remaining clauses prevent it from
changing the common equality-free grammar.
\index{defined-negation translation}
\end{hbdefinition}

\begin{hbtheorem}{thm:p07-lk-presentation-equivalence}{sketched=annex:p7.calculi.2}{Canonical/defined-negation LK equivalence}
\par\noindent
For every finite equality-free sequent \(\Gamma\Rightarrow\Delta\),
\[
 \Gamma\derives_{\mathsf{LK}}\Delta
 \quad\text{iff}\quad
 d(\Gamma)\derives_{\mathsf{LK}_{\mathrm d}}d(\Delta).
\]
The theorem compares existence of derivations only.  It asserts no identity
of derivation trees and no preservation of height, size, Cut occurrences, or
cut-freeness.
\end{hbtheorem}

\begin{hbdefinition}{def:p07-sequent-nd-translation}{Finite-succedent translation to classical natural deduction}
For a finite multiset \(\Delta\), let \(\bigvee\Delta\) be a fixed associated
disjunction of the distinct formulas in \(\operatorname{supp}(\Delta)\), in
the registry's formula order, and set
\(\bigvee\varnothing\defeq\bot\).  On the assumption side use the finite set
\(\operatorname{supp}(\Gamma)\).  This convention is required to translate
a multi-succedent sequent into the single-conclusion judgment of
\hyperref[sheet:nd-fol]{\texttt{sheet:nd-fol}}; it forgets multiplicity and
therefore is not a translation of proof objects.
\index{sequent!translation to natural deduction}
\end{hbdefinition}

\begin{hbtheorem}{thm:p07-lk-nd-equivalence}{cited=Gentzen1934a}{Derivability equivalence of LK and classical natural deduction}
For every finite equality-free sequent \(\Gamma\Rightarrow\Delta\), canonical
\(\mathsf{LK}\) derives \(\Gamma\Rightarrow\Delta\) iff the classical
natural-deduction calculus on
\hyperref[sheet:nd-fol]{\texttt{sheet:nd-fol}}, restricted to the
equality-free language, derives \(\bigvee\Delta\) from
\(\operatorname{supp}(\Gamma)\).  The propositional single-succedent case is
the cited Theorem~\ref{thm:p02-calculi-derivability-equivalence}.
\end{hbtheorem}

Starting instead with \(\mathsf{LK}_{\mathrm d}\) requires the translation
\(d\).  No equality version is inferred: it would have to name both
\texttt{sheet:lk-eq} and the exact natural-deduction equality rules, and it
would not inherit the cut-elimination result below.

\begingroup\scriptsize
\begin{longtable}{@{}p{.15\textwidth}p{.17\textwidth}p{.24\textwidth}
  p{.28\textwidth}@{}}
\caption{Calculi catalog and evidence boundaries (apparatus A-071).}
\label{tab:p07-calculi-catalog}\\
\toprule
Presentation & Registry home & Derivation object & Evidence recorded here \\
\midrule
\endfirsthead
\toprule
Presentation & Registry home & Derivation object & Evidence recorded here \\
\midrule
\endhead
Canonical \(\mathsf{LK}\) & \hyperref[sheet:lk]{\texttt{sheet:lk}} & Finite
tree of two-sided finite-multiset sequents & Propositional and first-order
soundness proved in Part~2; completeness sketched, with Choice in the
first-order form; no P7 cut-admissibility claim for this presentation. \\
\(\mathsf{LK}_{\mathrm d}\) &
\hyperref[sheet:lk-defneg]{\texttt{sheet:lk-defneg}} & Finite tree of
two-sided sequents with defined negation & Soundness and the translation to
canonical LK sketched here; Cut admissibility sketched below. \\
\(\mathsf{LK}_{=}\) & \hyperref[sheet:lk-eq]{\texttt{sheet:lk-eq}} &
Canonical LK tree with equality identity and replacement & Soundness proved
in Part~2; symmetric replacement uses Cut; no cut-admissibility or
completeness theorem for this equality extension is added here. \\
\(\mathsf{LJ}\) & \hyperref[sheet:lj]{\texttt{sheet:lj}} & Finite tree of
single-succedent sequents & Cut admissibility sketched below; semantic
adequacy belongs to the Part-8 IPL entry. \\
Classical natural deduction (the NK catalog alias) &
\hyperref[sheet:nd-cpl]{\texttt{sheet:nd-cpl}} and
\hyperref[sheet:nd-fol]{\texttt{sheet:nd-fol}} & Finite formula tree with
nonvacuous discharge, primitive negation, and RAA & Soundness and
completeness proved in Part~2; NK is only an alias for these exact sheets,
not a new active sheet. \\
Intuitionistic natural deduction \(\mathsf{NJ}\) &
\hyperref[sheet:nj]{\texttt{sheet:nj}} & Finite formula tree with the
sheet's exact discharge and eigenvariable conditions & Normalization sketched
in Chapter~\ref{sec:p07-normalization-and-curry-howard}; no semantic adequacy
result is manufactured here. \\
Hypersequent calculi & Citation-only survey; no active sheet & Finite trees of
multi-component sequents, presentation dependent & No adequacy,
cut-admissibility, or decision result is assigned without a selected sheet. \\
Negri's labelled modal sequent calculi & Citation-only survey; no active sheet
& Finite trees whose sequents contain formulas annotated by world labels and
syntactic relational atoms; unlike ordinary sequents, the selected frame
relation is itself syntactic data & \cite{Negri2005} is the source
pointer; no exact P7 sheet, derivation, or adequacy claim is assigned. \\
Display calculi & Citation-only survey; no active sheet & Structured
consecution trees, presentation dependent & No local evidence claim; this row
is a forward pointer only. \\
Deep-inference calculi & Citation-only survey; no active sheet & Derivations
whose inferences may act inside formula or structure contexts & No local
evidence claim; exact syntax and congruence must be fixed first. \\
\bottomrule
\end{longtable}
\endgroup

\paragraph{Source disposition.}
The v1.1 rule displays have moved verbatim in substance to
\texttt{sheet:lk-defneg} and \texttt{sheet:lj}; the body does not duplicate
them.  Its worked cut-free derivation of
\(\varphi\land\psi\Rightarrow\psi\land\varphi\) is retired because the
existing worked comparison in Figure~\ref{fig:p02-nd-lk-correspondence}
supplies its lookup function without a second proof display.  The source
exercise requesting an LJ distributivity derivation is retired with its
principal-formula lesson retained in
Definition~\ref{def:p07-principal-formula}; the eigenvariable-unsoundness
exercise is retired with its lesson preserved as the binding freshness
conditions on the active sheets.  Nothing in these dispositions is relabeled
as a proof.
